% !TeX root = ../main.tex
\section{Case Study} %
\label{sec:case}

We illustrate the expressiveness of context types with three
case studies: compositional may- and must-style reasoning
(\autoref{subsec:may-must}), a coverage completeness property for a BST
generator (\autoref{subsec:coverage}), and disprove reasoning for a merge
function (\autoref{subsec:disprove}).

\subsection{May and Must Analysis}\label{subsec:may-must}

\begin{figure}[h!]
    \centering
    \begin{subfigure}[t]{0.48\linewidth}
        \centering
\begin{minted}[xleftmargin=5pt, numbersep=4pt, linenos = true, fontsize = \footnotesize, escapeinside=??]{OCaml}
val g : ?$\Code{i}{:}\ort{\Int}{\top}\ctsarr\ort{\Int}{0 \leq \vnu}$?

?$\vdash$? let f (x: int) (y: int) =
    (g x) < 0 || (g y) < 0
: ?$\Code{x}{:}\ort{\Int}{\top}\ctsarr \Code{y}{:}\ort{\Int}{\top}\ctsarr \ort{\Bool}{\vnu = \false}$?
\end{minted}
        \subcaption{May-analysis helps may-analysis.}
        \label{fig:may-must-may}
    \end{subfigure}
    \hfill
    \begin{subfigure}[t]{0.48\linewidth}
        \centering
\begin{minted}[xleftmargin=5pt, numbersep=4pt, linenos = true, fontsize = \footnotesize, escapeinside=??]{OCaml}
val g : ?$\Code{i}{:}\urt{\Int}{\top}\ctsarr\urt{\Int}{0 \leq \vnu}$?

?$\vdash$? let f (x: int) (y: int) =
    0 <= (g x) && 0 <= (g y)
: ?$\Code{x}{:}\urt{\Int}{\top}\ctsarr\Code{y}{:}\urt{\Int}{\top}\ctsarr\urt{\Bool}{\vnu = \true}$?
\end{minted}
        \subcaption{Must-analysis helps may-analysis.}
        \label{fig:may-must-must}
    \end{subfigure}

    \medskip

    \begin{subfigure}[t]{0.48\linewidth}
        \centering
\begin{minted}[xleftmargin=5pt, numbersep=4pt, linenos = true, fontsize = \footnotesize, escapeinside=??]{OCaml}
val g : ?$\Code{i}{:}\ort{\Int}{\top}\ctsarr\ort{\Int}{\vnu \neq 0}$?

?$\vdash$? let f (x: int) (y: int) =
    if y > 10 then (g x) + 1
    else x + 1
: ?$\Code{x}{:}\urt{\Int}{\top}\ctsarr\Code{y}{:}\urt{\Int}{\top}\ctsarr\urt{\Int}{\vnu = 1}$?
\end{minted}
        \subcaption{May-analysis helps must-analysis.}
        \label{fig:may-must-reach2}
    \end{subfigure}
    \hfill
    \begin{subfigure}[t]{0.48\linewidth}
        \centering
\begin{minted}[xleftmargin=5pt, numbersep=4pt, linenos = true, fontsize = \footnotesize, escapeinside=??]{OCaml}
val g : ?$\Code{i}{:}\ort{\Int}{\top}\ctsarr\ort{\Int}{\vnu > 0}\interty $?
        ?$\Code{i}{:}\urt{\Int}{\top}\ctsarr\urt{\Int}{\vnu > 0}$?

?$\vdash$? let f (x: int) (y: int) =
    let a = g x in
    if a > 10 then 3
    else a
: ?$\Code{x}{:}\ort{\Int}{\top}\ctsarr\Code{y}{:}\ort{\Int}{\top}\ctsarr\ort{\Int}{0 \leq \vnu}$?
\end{minted}
        \subcaption{Must-analysis helps may-analysis.}
        \label{fig:may-must-inter}
    \end{subfigure}
    \caption{Compositional may- and must-style reasoning with context types.}
    \label{fig:may-must-analysis}
\end{figure}

Our unified treatment of safety and reachability through context types
also furnishes a foundation for compositional may- and must-style
program analysis, as developed in prior work~\cite{GN+10}. That
analysis is compositional because callee summaries inform caller
analysis; Figure~\ref{fig:may-must-analysis} reproduces four
representative examples from the original paper. We restate them in
our core language and assign context types to the caller $\Code{f}$
and the callee $\Code{g}$. In our encoding, may-analysis summaries use
demonic parameters and returns, whereas must-analysis summaries use
angelic parameters and returns.

Figure~\ref{fig:may-must-analysis}~(a) calls for a ``not-may'' summary
of $\Code{f}$, namely a safety property stating that $\Code{f}$'s
result \emph{may not} be $\true$. Because the callee $\Code{g}$
returns only non-negative values (i.e., $\ort{\Int}{0 \leq \vnu}$),
the caller $\Code{f}$ may be summarized as returning only $\false$. By
contrast, the must-analysis in Figure~\ref{fig:may-must-analysis}~(b)
requires $\true$ to be reachable for $\Code{f}$, relying on
$\Code{g}$'s guarantee that every non-negative value is
reachable. Prior must-style analyses can handle only plain
conjunctions of two callee return values, since they cannot express
disjointness between the inputs $\Code{x}$ and $\Code{y}$ in
$\Code{f}$. On the other hand, if function $\Code{f}$ returns
$\Code{(g\;x) - (g\;y)}$ and we want to prove all natural number are
reachable, the following disjointed type for $\Code{f}$ is required:
\begin{align*}
  \Code{x}{:}\urt{\Int}{\top}\ctwand\Code{y}{:}\urt{\Int}{\top}\ctwand\urt{\Int}{0 \leq \vnu }
\end{align*}\noindent

Figure~\ref{fig:may-must-analysis}~(c) rules out reachability analysis
for the then-branch of the conditional, because $\Code{g}$ may return
a non-zero value and thereby prevent that branch from reaching the
result $1$. In Figure~\ref{fig:may-must-analysis}~(d), by contrast,
the callee $\Code{g}$ carries two summaries as an intersection type;
its angelic component states that every integer is reachable. This
forces may-analysis of the caller $\Code{f}$ to explore both branches
of the conditional. Our type system handles such cases more uniformly,
since a single type may be constrained under both angelic and demonic
modalities. For instance, safety checking for $\Code{f}$ can proceed
under the more comprehensive type:
\begin{align*}
    \Code{x}{:}(\ort{\Int}{\top}\interty \urt{\Int}{\top})\ctsarr\Code{y}{:}(\ort{\Int}{\top}\interty \urt{\Int}{\top})\ctsarr(\ort{\Int}{\top}\interty \urt{\Int}{\vnu = 1})
\end{align*}\noindent
which supports safety and reachability checking in one pass, rather than alternating between may- and must-style analysis.

\subsection{Coverage Completeness Verification}\label{subsec:coverage}
The coverage type system of~\citet{CoverageType} verifies that a
test generator produces every desirable input.
Consider the following BST generator:
\begin{minted}[xleftmargin=15pt, numbersep=4pt, linenos = true, fontsize = \footnotesize, escapeinside=??]{OCaml}
let rec bst_gen (r: nat) (lo: nat)
              (range: nat -> nat -> nat): nat tree =
if r <= 0 then Leaf
else Leaf ?$\oplus$?
  let r' = range (1, r) in
  let left  = bst_gen (r' - 1) lo range in
  let right = bst_gen (r - r') (lo + r') range in
  Node (lo + r') left right
\end{minted}
The following specification asserts that for any valid range generator,
\texttt{bst\_gen} produces every BST whose keys fall in $(\Code{lo},
\Code{lo}+r]$:
\begin{align*}
  \mykw{val}\;\Code{bst\_gen} : \ctpers(&
    \Code{r}{:}\ort{\Nat}{\top} \ctsarr \Code{lo}{:}\ort{\Nat}{\top} \ctsarr
  \\&
    \Code{range}{:}\ctpers (x{:}\ort{\Nat}{\top} \ctsarr
      y{:}\ort{\Nat}{x < \vnu} \ctsarr
      \urt{\Nat}{x < \vnu < y})
    \ctsarr
  \\&
    \urt{\Nat\;\Code{tree}}{\Code{bst}(\vnu) \land
      \forall x.\, \Code{mem}(\vnu, x) \iff
      (\Code{lo} < x \leq \Code{lo}+r)})
\end{align*}\noindent
In our setting, all such generators can \emph{produce new angelic capabilities} for \emph{arbitrary} input parameters and thus should have context types of the following form:
\begin{align*}
    \overbrace{\ctpers (\quad \underbrace{\overline{x{:}\ort{b_x}{\phi_x}\ctsarr}}_{\texttt{demonic parameters}} \quad \underbrace{\urt{b}{\phi}}_{\texttt{angelic result}} \quad )}^{\texttt{duplicatable}}
\end{align*}\noindent
Thus, the range generator is persistent and has demonic parameter types ($\ort{\Nat}{\top}$ and
$\ort{\Nat}{x < \vnu}$) and an angelic result type ($\urt{\Nat}{x <
  \vnu < y}$).  The demonic parameters mean that successive calls to
\texttt{range} with different arguments can be typed with
independent result capabilities — the generator's output at each
recursive call is disjoint from its outputs at other calls.
Similarly, the BST generator follows the same pattern: its demonic parameter types allow any valid size $\Code{r}$ and lower bound $\Code{lo}$, and its angelic result type asserts that every valid BST is reachable by some sequence of generator choices.

\subsection{Disprove Reasoning}\label{subsec:disprove}


\begin{figure}[h!]
    \centering
\begin{minted}[xleftmargin=5pt, numbersep=4pt, linenos = true, fontsize = \footnotesize, escapeinside=??]{OCaml}
    let rec merge (l1: elem list) (l2: elem list) : elem list =
        match l1, l2 with
        | l1, [] -> l1
        | [], l2 -> l2
        | h1 :: t1, h2 :: t2 ->
            let t = merge t1 t2 in (* t contains elements in t1 *)
            if h1 < h2 then h1 :: h2 :: t (* h2 may be greater than elements in t1 *)
            else if h1 > h2 then h2 :: h1 :: t
            else h1 :: t
\end{minted}
\caption{A function that merges two sorted lists. If either of
the input lists are empty, \Code{merge} returns the other (lines
$3-4$). Otherwise, the function merges the tails of the two inputs and appends the two top elements ahead of the merged result. }
\label{fig:merge}
\end{figure}

When a refinement type system fails to prove a safety property, this does not mean the program is incorrect: may-style reasoning offers no guarantee that erroneous behaviors are reachable. Recent work~\cite{RefinementTypeRefutations,RefinementTypeSymbolicExecution} addresses this limitation by using symbolic execution to search for counterexamples; our context type system can establish such counterexamples directly.
Consider \Code{merge} in Figure~\ref{fig:merge}, a core component of merge sort that should preserve sortedness of the output list:
\begin{align*}
  \Code{l1}{:}\ort{\List[\Code{elem}]}{\top} \ctsarr \Code{l2}{:}\ort{\List[\Code{elem}]}{\top} \ctsarr \ort{\List[\Code{elem}]}{\I{sorted}(\vnu)}
\end{align*}\noindent
Verification fails nonetheless. The base cases (lines $3$--$4$) type-check, so the failure arises in the recursive case, where the refinement type system offers little guidance.
Symbolic execution can find the following counterexample:
\begin{align*}
    \texttt{Counterexample:}&\quad  \text{ when } \Code{l1} = [1;2] \text{ and } \Code{l2} = [3;4] \text{ then } \mstep{\Code{merge}\;[1;2]\;[3;4]}{[1;3;2;4]}
\end{align*}\noindent
The bug stems from the assumption that, after comparing $\Code{h1}$ and $\Code{h2}$, $\Code{h2}$ is always the second-smallest element (line $7$).
Even when $\Code{h1} < \Code{h2}$, $\Code{h2}$ may exceed some element of $\Code{t1}$, breaking sortedness.
Alternatively, we can \emph{disprove} the function with the following type:
\begin{align*}
    &\Code{l1}{:}\urt{\List[\Code{elem}]}{\top} \ctwand \Code{l2}{:}\urt{\List[\Code{elem}]}{\top} \ctwand \ort{\List[\Code{elem}]}{\neg \I{sorted}(\vnu)} \ctoplus \ort{\List[\Code{elem}]}{\top}
\end{align*}\noindent
This type follows the ``affine reachability'' style of \autoref{sec:context-typing}: the return type is a sum of the top type $\ort{\List[\Code{elem}]}{\top}$ and the negation of the original safety property, $\ort{\List[\Code{elem}]}{\neg \I{sorted}(\vnu)}$, expressing erroneous behavior.
The parameter types are moreover flipped to \emph{disjointed angelic types}, indicating that $\Code{l1}$ and $\Code{l2}$ may independently range over arbitrary sorted lists.
Without this disjointness, a carefully crafted environment can hide the bug:
\begin{center}
\begin{minted}[xleftmargin=5pt, numbersep=4pt, linenos = true, fontsize = \footnotesize, escapeinside=??]{OCaml}
  let l1 = sorted_list_gen () in
  let l2 = match l1 with
           | [] -> sorted_list_gen ()
           | _ -> [] in
  merge l1 l2
\end{minted}
\end{center}\noindent
Here $\Code{l1}$ and $\Code{l2}$ each range over all sorted lists, yet they cannot both be non-empty simultaneously.
Under this environment, the merge bug cannot be triggered, since the case where both inputs are non-empty (line $5$) is never reached.
